\documentclass[../main.tex]{subfiles}

\begin{document}

\section{INTRODUCTION}


\subsection{Background}
Poultry farming is a key industry of agriculture that contributes significantly towards enhancing food security and economic development across the world. In Nepal, poultry farming produces approximately 1.61 billion eggs and 548,000 tons of poultry meat annually. The industry not only provides sufficient products to meet the increasing demand for protein food but also provides employment for over 55,000 individuals in commercial chicken rearing.

Maintaining ideal environmental conditions such as temperature, humidity, ventilation, and gas composition inside poultry farms is, however, a major challenge. Suboptimal conditions can lead to high mortality, reduced growth, and low production, which further affect farm profitability.

The Internet of Things (IoT) is now a transformative technology in agriculture, making real-time tracking and automation of various processes possible. With sensors, microcontrollers, and cloud computing, IoT devices provide farmers with actionable insights and visibility into their operations, enhancing efficiency and reducing losses. The project leverages IoT technology to address these challenges, intending to build a smart, data-heavy system for managing poultry farms.
\subsection{Motivation}
  The traditional practices for monitoring poultry farms rely on manual observation, which are labor-intensive, time-consuming, and prone to errors. These practices often produce delayed feedback in case of adverse conditions, negatively impacting bird health and productivity. In addition, the lack of real-time data hinders data-driven decision-making and optimized operations.\\
An IoT-based system has the potential to provide solutions to these issues through continuous monitoring and automated control of environmental parameters. Such a system can promptly alert farmers to deviations from optimal conditions, enabling timely interventions. Automation also reduces the need for manual labor, thereby lowering operational costs and increasing efficiency.
\subsection{Statement of the Problem}
The primary problem addressed by this project is the absence of real-time monitoring and automation in poultry farms, leading to suboptimal environmental conditions and increased operational costs. Specific challenges include:
\begin{itemize}
\item Unsteady temperature and humidity that affect chicken health and growth.

\item Elevated levels of ammonia and CO₂ gases that are harmful to poultry health.

\item Inaccurate and labor-intensive manual monitoring practices.

\item Difficulty in observing multiple parameters at once.
\end{itemize}
\subsection{Project objective}
The objectives of this project are:
\begin{itemize}
    \item To install a system for real-time monitoring of significant environmental control parameters of poultry farms, i.e., temperature, humidity, ammonia, CO₂, and chicken weight.
   \item To implement automated control systems to provide the best conditions, i.e., controlling ventilation and heat.
   \item To create a simple-to-use web or mobile dashboard for farmers to see real-time data, receive notifications, and control the system remotely.
   % \item To evaluate the system’s performance and reliability through testing and validation.
\end{itemize}
\subsection{Significance of the study}
This project aims to transform traditional poultry farming practices into smart, data-driven operations using IoT technology. Ensuring the best environmental conditions can enhance animal welfare, increase productivity, and reduce the cost of operations for farmers. It becomes an ideal case study for integrating technology in agricultural areas, promoting sustainable and efficient farming practices in Nepal and across the world.

The successful implementation of the project would also serve as a model for applying similar technologies in other  sectors, in pursuit of the larger goal of streamlining agriculture.
\end{document}